\documentclass[11pt]{article}
\usepackage{amsmath,amssymb,amsthm}
\usepackage{booktabs}
\usepackage{enumitem}
\usepackage{hyperref}
\usepackage[margin=1in]{geometry}

\hypersetup{colorlinks=true,linkcolor=black,citecolor=blue,urlcolor=blue}

\newtheorem{theorem}{Theorem}
\newtheorem{lemma}[theorem]{Lemma}
\newtheorem{corollary}[theorem]{Corollary}
\newtheorem{definition}{Definition}
\newtheorem{remark}{Remark}
\newtheorem{proposition}[theorem]{Proposition}

\title{Zoo DAO Governance Soundness\\
{\large Homomorphic / Holographic-Consensus Voting with FHE-Tallied
Secrecy and Sybil Bounds}}
\author{Zoo Labs Foundation\\
\texttt{research@zoo.ngo}}
\date{Spec freeze: 2026-04-15 \\ Effective: 2026-04-20}

\begin{document}
\maketitle

\begin{abstract}
The Zoo DAO uses a holographic-consensus governance scheme with
weighted voting where weights are derived from advocacy /
involvement / contribution scores stored on-chain. Vote tallies are
computed under Fully Homomorphic Encryption (TFHE / CKKS, see
\cite{lp-fhe}) so that individual votes remain secret while the
total is publicly verifiable. We prove three results:
(i)~weighted-voting math correctness under the homomorphic tally;
(ii)~vote-secrecy preservation against a passive adversary holding
the chain transcript;
(iii)~an anti-Sybil bound on advocacy / involvement / contribution
(AIC) scores.
\end{abstract}

\tableofcontents

%% ==================================================================
\section{Background}
\label{sec:bg}
%% ==================================================================

The Zoo DAO governance contract implements weighted voting on
proposals. Each member $m$ has an AIC score
$s_m = \alpha \cdot a_m + \beta \cdot i_m + \gamma \cdot c_m$
where $a_m$ is advocacy (community participation), $i_m$ is
involvement (engagement metrics), $c_m$ is contribution (verified
on-chain commits / payments / NGO actions). Weights
$\alpha, \beta, \gamma$ are chain-genesis constants.

Votes on proposal $P$ are submitted as ciphertexts
$c_m = \mathrm{Enc}_{\mathrm{tfhe}}(v_m)$ where $v_m \in \{0, 1\}$
(reject / approve). The DAO contract aggregates ciphertexts using
TFHE addition; the total is decrypted under threshold decryption
(M-Chain MPC, LP-134) and compared against a quorum threshold.

%% ==================================================================
\section{Weighted-voting correctness}
\label{sec:weight}
%% ==================================================================

\begin{definition}[Weighted vote ciphertext]
\label{def:wv}
A weighted vote is
$c_m^\star = \mathrm{Enc}_{\mathrm{tfhe}}(s_m \cdot v_m)$, where
$s_m$ is public (stored on-chain) and $v_m$ is the secret choice.
\end{definition}

\begin{lemma}[Linearity of TFHE addition]
\label{lem:linear}
For TFHE ciphertexts $c_1, c_2$ encrypting integers $x_1, x_2$ in
the homomorphic plaintext space,
$\mathrm{Dec}(c_1 \boxplus c_2) = x_1 + x_2$.
\end{lemma}

\begin{proof}
Standard property of TFHE; see \cite{tfhe} Theorem 4.1.
\end{proof}

\begin{theorem}[Tally correctness]
\label{thm:tally}
Given weighted vote ciphertexts
$\{c_m^\star\}_{m \in \mathcal{M}}$, the tally
$T = \sum_{m \in \mathcal{M}} s_m \cdot v_m$ computed as
$\mathrm{Dec}(\boxplus_m c_m^\star)$ is correct: the decrypted
result equals the plaintext weighted sum.
\end{theorem}

\begin{proof}
By Lemma~\ref{lem:linear} and induction on the number of
ciphertexts. The threshold decryption step is correct under the
M-Chain MPC unforgeability assumption (see \cite{cggmp21-lean,
frost-lean}).
\end{proof}

\begin{remark}[Quorum and approval]
A proposal $P$ passes iff $T \geq q \cdot \sum_m s_m$ for
quorum constant $q$. Both $T$ and $\sum_m s_m$ are publicly known
post-tally; the comparison is a public deterministic predicate.
\end{remark}

%% ==================================================================
\section{Vote secrecy preservation}
\label{sec:secrecy}
%% ==================================================================

\begin{definition}[Secrecy game]
\label{def:secrecy}
The vote-secrecy game between adversary $\mathcal{A}$ and the
DAO is:
\begin{enumerate}[nosep]
  \item $\mathcal{A}$ chooses target member $m^\star$;
  \item the DAO sets $v_{m^\star} \overset{\$}{\leftarrow} \{0,1\}$
        and publishes $c_{m^\star}^\star =
        \mathrm{Enc}_{\mathrm{tfhe}}(s_{m^\star} \cdot v_{m^\star})$;
  \item $\mathcal{A}$ outputs guess $\tilde{v}$.
\end{enumerate}
$\mathcal{A}$ wins iff $\tilde{v} = v_{m^\star}$.
$\mathsf{Adv}^{\mathrm{secrecy}}_{\mathcal{A}} =
|\Pr[\mathcal{A}\text{ wins}] - 1/2|$.
\end{definition}

\begin{theorem}[Vote secrecy]
\label{thm:secrecy}
Under TFHE IND-CPA security
(\cite{tfhe} Theorem 5.2),
$\mathsf{Adv}^{\mathrm{secrecy}}_{\mathcal{A}} \leq
\mathsf{Adv}^{\mathrm{IND\text{-}CPA}}_{\mathrm{TFHE}} =
\mathsf{negl}(\lambda)$.
\end{theorem}

\begin{proof}
Reduction: any distinguisher of $v_{m^\star} = 0$ vs $v_{m^\star}=1$
ciphertexts gives an IND-CPA distinguisher on TFHE. The
threshold-decrypted tally only reveals the sum, not individual
votes; the sum information leak is bounded by entropy of
$|\mathcal{M}|$-fold sums (a standard mix-net-style argument).
\end{proof}

\begin{corollary}[Coercion resistance under threshold decryption]
\label{cor:coerc}
A coercer who observes only on-chain ciphertexts cannot learn the
target member's vote with non-negligible advantage. Coercion
resistance against the threshold-decrypter committee is bounded by
the M-Chain MPC committee's honesty: at least one of the
$t$ decrypters must be honest for vote secrecy to hold, where
$t$ is the threshold parameter.
\end{corollary}

%% ==================================================================
\section{Anti-Sybil bound on AIC scores}
\label{sec:sybil}
%% ==================================================================

\begin{definition}[AIC score growth]
\label{def:aic}
The AIC score $s_m$ for member $m$ at time $t$ is
\[
  s_m(t) = \alpha \int_0^t a_m(\tau) d\tau +
            \beta \int_0^t i_m(\tau) d\tau +
            \gamma \cdot \mathrm{Contrib}_m(t),
\]
where $a_m, i_m$ are continuous-time non-negative engagement signals
(rate-limited by the DAO's anti-spam policy) and
$\mathrm{Contrib}_m(t)$ is the cumulative on-chain contribution
verified under triple-hybrid identity (Crypto/Hybrid.lean).
\end{definition}

\begin{lemma}[Engagement rate bound]
\label{lem:engage}
The DAO contract enforces
$a_m(\tau), i_m(\tau) \leq R_{\max}$ for some rate
constant $R_{\max}$ via on-chain rate-limit checks. Hence
$\alpha a_m + \beta i_m \leq (\alpha + \beta) R_{\max}$ at all
times.
\end{lemma}

\begin{theorem}[Sybil bound]
\label{thm:sybil}
A Sybil set $\mathcal{S}$ of $k$ fake identities, each generating
the maximum allowed engagement and zero verified contribution,
contributes at most $k \cdot (\alpha + \beta) R_{\max} \cdot t$
to the total AIC score after time $t$. Honest contribution-weighted
votes from a single real identity with verified
$\mathrm{Contrib}_m(t) > k \cdot (\alpha + \beta) R_{\max} \cdot t /
\gamma$ outvotes the Sybil set.
\end{theorem}

\begin{proof}
By Lemma~\ref{lem:engage}, the Sybil aggregate AIC growth is
bounded. Verified contribution $\mathrm{Contrib}_m$ requires
on-chain transactions or NGO partnership records, which Sybils
cannot fabricate at zero cost. Choosing
$\gamma$ large relative to $(\alpha + \beta) R_{\max}$ in the
chain genesis closes the loop; the explicit threshold is
$\mathrm{Contrib}_m > (k (\alpha + \beta) R_{\max} t)/\gamma$.
\end{proof}

\begin{remark}[Practical parameters]
The 2026 Zoo DAO uses $\alpha = 1, \beta = 1, \gamma = 100$ and
$R_{\max} = 100$ engagement-units/day. A single $\$10$ verified
contribution thus exceeds 10 days of Sybil-maxed engagement from
one fake account; 100-fake-account Sybils are bounded by 1 verified
contribution of $\$100$.
\end{remark}

%% ==================================================================
\section{Detailed reduction: vote secrecy}
\label{sec:detail-secrecy}
%% ==================================================================

We expand the proof of Theorem~\ref{thm:secrecy}.

\paragraph{Game.}
Following the standard IND-CPA-style framing for TFHE:
\begin{enumerate}[nosep]
  \item $\mathcal{A}$ submits the index $m^\star$ of the target
        member; the simulator commits to all other members'
        votes (chosen by $\mathcal{A}$, modeling honest-but-known
        counters).
  \item The simulator picks $b \overset{\$}{\leftarrow} \{0,1\}$
        and sets $v_{m^\star} = b$.
  \item The simulator computes
        $c_{m^\star}^\star = \mathrm{Enc}_{\mathrm{tfhe}}(s_{m^\star} \cdot b)$.
  \item $\mathcal{A}$ outputs guess $\tilde{b}$.
\end{enumerate}
$\mathcal{A}$'s advantage is $|\Pr[\tilde{b} = b] - 1/2|$.

\paragraph{Reduction to TFHE IND-CPA.}
Given $\mathcal{A}$ in the secrecy game, build a TFHE distinguisher
$\mathcal{D}$:
\begin{enumerate}[nosep]
  \item $\mathcal{D}$ receives a TFHE challenge ciphertext $c^\star$
        encrypting either $0$ or $s_{m^\star}$ (the simulator
        chose).
  \item $\mathcal{D}$ runs $\mathcal{A}$ on input $c^\star$; outputs
        $\mathcal{A}$'s guess as the distinguishing bit.
\end{enumerate}
$\mathsf{Adv}_{\mathcal{D}}^{\mathrm{IND\text{-}CPA}} =
 \mathsf{Adv}_{\mathcal{A}}^{\mathrm{secrecy}}$ (modulo the trivial
sum-information, which has zero contribution under the standard
information-theoretic mix-net argument because the sum is
decrypted only after all ciphertexts are committed).

\paragraph{Conclusion.}
$\mathsf{Adv}^{\mathrm{secrecy}} \leq
 \mathsf{Adv}^{\mathrm{IND\text{-}CPA}}_{\mathrm{TFHE}}$, the
right-hand side negligible by~\cite{tfhe} Theorem 5.2.

%% ==================================================================
\section{Worked example: a typical proposal}
\label{sec:example}
%% ==================================================================

Proposal $P$: ``Disburse 50{,}000 USDL to WWF for 2026 conservation
program.'' Members $m_1, \ldots, m_5$ have AIC scores
$s_1 = 100, s_2 = 200, s_3 = 50, s_4 = 30, s_5 = 70$.
Members $\{m_1, m_2, m_5\}$ vote approve ($v=1$); $\{m_3, m_4\}$
vote reject ($v=0$).

Tally:
$T = 100 \cdot 1 + 200 \cdot 1 + 50 \cdot 0 + 30 \cdot 0 + 70 \cdot 1 = 370$.

Total stake $\sum s = 450$; quorum $q = 0.5$ requires $T \geq 225$.
$370 \geq 225$, proposal passes.

By Theorem~\ref{thm:tally}, the homomorphic sum $T$ is correct.
By Theorem~\ref{thm:secrecy}, an adversary observing only the
on-chain ciphertexts learns nothing about individual votes.
Sybil bound (Theorem~\ref{thm:sybil}) ensures $m_3, m_4$ cannot
multiply themselves into a winning coalition without
observable on-chain contributions.

%% ==================================================================
\section{Threshold-decryption ceremony}
\label{sec:thresh-dec}
%% ==================================================================

The decryption is a ceremony on M-Chain (LP-134):

\begin{enumerate}[nosep]
  \item DAO contract closes voting at deadline; the encrypted
        sum $C_T = \boxplus_m c_m^\star$ is published in the DAO
        contract storage.
  \item M-Chain validators run the FROST-style threshold decryption
        ceremony for $C_T$, requiring $t$-of-$n$ honest decrypters
        (\cite{frost-lean}).
  \item The plaintext $T$ is published; the DAO contract compares
        against $q \cdot \sum_m s_m$ and emits the
        \texttt{ProposalPassed} or \texttt{ProposalFailed} event.
\end{enumerate}

The decryption ceremony itself is finalised on M-Chain under
Quasar 3.0 (\cite{quasar-cert}); the result is therefore
PQ-final.

%% ==================================================================
\section{Cross-references}
\label{sec:cross}
%% ==================================================================

\begin{itemize}[nosep]
  \item ZIP-0017 -- DAO governance core.
  \item ZIP-0023 -- Holographic-consensus reputation.
  \item ZIP-0026 -- Quadratic / weighted voting framework.
  \item ZIP-0027 -- AIC score derivation.
\end{itemize}

%% ==================================================================
\section{Composition with Liquidity Protocol}
\label{sec:lp}
%% ==================================================================

DAO treasury disbursements (e.g.\ to NGOs) settle via the Liquidity
Protocol post 2026-04-20 (\cite{zoo-adopts}). Composition is direct:
the DAO emits a settlement intent $(t, P, \mathrm{NGO}_i, \mathrm{amount})$
when a passing tally is decoded; the OMA contract on Lux C-Chain
executes the swap. Atomicity follows
\cite[Theorem 5.2]{liquidity-lp}; PQ finality follows
\cite[Theorem 7.7]{quasar-cert}.

%% ==================================================================
\begin{thebibliography}{9}

\bibitem{tfhe}
I.~Chillotti, N.~Gama, M.~Georgieva, M.~Izabach\`ene.
\newblock TFHE: Fast Fully Homomorphic Encryption over the Torus.
\newblock J.~Cryptology, 2020.

\bibitem{lp-fhe}
Lux Industries.
\newblock LP-FHE Series (CKKS, TFHE).
\newblock \texttt{luxfi/lps/}.

\bibitem{cggmp21-lean}
Lux Industries.
\newblock \texttt{lean/Crypto/Threshold/CGGMP21.lean}.

\bibitem{frost-lean}
Lux Industries.
\newblock \texttt{lean/Crypto/FROST.lean}.

\bibitem{quasar-cert}
Z.~Kelling.
\newblock QuasarCert Soundness, Quasar 3.0.
\newblock \texttt{luxfi/proofs/quasar-cert-soundness.tex}, 2025-12-15.

\bibitem{liquidity-lp}
Liquidity.io.
\newblock The Liquidity Protocol v1.0.
\newblock \texttt{liquidityio/proofs/liquidity-protocol/}, April 2026.

\bibitem{zoo-adopts}
Zoo Labs Foundation.
\newblock Zoo Adopts the Liquidity Protocol.
\newblock \texttt{zooai/proofs/zoo-adopts-liquidity-protocol.tex}, 2026-04-20.

\end{thebibliography}

\end{document}
